\newpage
\section{More Related Works}
\label{sec:more_related_works}
\paragraph{Multimodal Large Language Models. } Multimodal Large Language Models (MLLMs) take inputs beyond text (\textit{e.g.}, image, audio, video, \textit{etc}) and generate text responses \cite{huang2023language}. Most MLLMs focus on vision-language tasks. Prototypical approaches train adaptors that connect independent visual-only and language-only modules \cite{li2022blip, li2023blip, alayrac2022flamingo} or adapt language models to visual inputs \cite{huang2023language, chen2022pali, chen2023pali}. With instruction tuning \cite{wei2021finetuned} and accessibility to more instruction-tuned Large Language Models \cite{touvron2023llama, jiang2023mistral, young2024yi, vicuna2023}, there has been a proliferation of open-source MLLMs \cite{liu2023llava, zhang2023llama, zhu2023minigpt, ye2023mplug, dai2023instructblip, li2023mimic, laurençon2023obelics, lin2023sphinx, fuyu-8b}. More recent work has attempted to scale up the backbone language model, add more alignment data, increase input resolution, design different vision-language adaptation paradigms, and finetune more modules that are otherwise frozen to improve the capabilities of MLLMs \cite{liu2023improvedllava, liu2024llavanext, chen2023internvl, dong2024internlmxcomposer24khd, li2024minigemini, chen2024far, laurenccon2024matters, gao2024sphinxx, lee2024moai, lu2024deepseekvl}. While many recent open-source MLLMs reported on-par or better performance compared to proprietary models in chart understanding \cite{lu2024mathvista, masry2022chartqa}, little is known about how well these models generalize. In our work, we evaluate the most recent MLLMs on modified versions of chart subsets from MathVista \cite{lu2024mathvista} (\cref{sec:motivation}) and \ours{} (\cref{sec:experiments}), showing that open-source models generalize poorly and the performance gap still exists.

\paragraph{MLLM Benchmarks. } Prototypical MLLM benchmarks follow Visual Question Answering based on natural images \cite{antol2015vqa, goyal2017making, hudson2019gqa, singh2019towards, marino2019ok} or \textit{screenshots} \cite{gao2024improving}, such as documents \cite{mathew2021docvqa}, diagrams \cite{kembhavi2016diagram}, charts \cite{masry2022chartqa} and infographics \cite{mathew2022infographicvqa}. More recently, several MLLM benchmarks emerged that evaluate multimodal capabilities in a more \emph{knowledge-intensive} \cite{lu2022learn, lu2024mathvista, yue2023mmmu} and \emph{comprehensive} \cite{yu2023mm, liu2024mmbench, fu2024mme} setting. Chart understanding signifies an important challenge for MLLMs, where the vast majority of open- and proprietary models \cite{openai2024gpt4, claude3, rekateam2024reka, bai2023qwenvl, team2023gemini} report model performance on chart understanding tasks \cite{lu2024mathvista, masry2022chartqa}. Earliest chart understanding benchmarks often adopt synthetic data and charts \cite{kahou2017figureqa, kafle2018dvqa, methani2020plotqa} or use stylistically consistent charts \cite{masry2022chartqa}. More recent chart understanding benchmarks are either not publicly available \cite{liu2024mmc, xu2024chartbench} or widely adopted \cite{xia2024chartx}. \ours{} (\cref{sec:charxiv_overview}) is most similar to the design choice of ChartQA \cite{masry2022chartqa}, yet we adopt more natural, diverse and challenging charts with human-curated QA pairs, resulting in a benchmark that better reflects general capabilities in chart understanding.